\chapter{Conclusion}
\label{chap:conclusion}

\section{Summary of Contributions}

This thesis addressed the problem of evaluating credit-default prediction
beyond a single aggregate performance metric. It combined explainability,
age-group fairness, proxy auditing, and mitigation in one experimental design
so that each analysis could be interpreted against the same data split and
decision threshold.

The main methodological contribution is the fairness-oriented audit of data
preparation. The earlier pipeline retained duplicates, removed older
applicants through distributional outlier filtering, and added ratio features
that recovered age exactly. The revised pipeline removes only logically
invalid records, keeps missingness, avoids age-derived features, and stratifies
on age group and outcome.

The preliminary baseline comparison selected LightGBM, which achieved ROC-AUC
0.9431 under the earlier protocol. The final claims about SHAP agreement,
fairness gaps, residual proxy information, and reweighing cost must be based
on the final experiment outputs. This distinction is important: a model can
rank risk well while still producing unequal approval behavior.

\section{Revisiting the Research Objectives}

\begin{enumerate}
  \item \textbf{Develop a credit-default classifier.} Achieved at the
  methodological level with LightGBM and a fixed, reproducible configuration.
  The preliminary baseline achieved ROC-AUC 0.9431.
  \item \textbf{Explain model decisions.} The evaluation protocol defines
  gain, split, and TreeSHAP importance and their rank-agreement measures.
  Final ranking values must be inserted after the experiment is run.
  \item \textbf{Measure age-group fairness.} The study defines DPD, EOD,
  false-positive-rate differences, group-wise ROC-AUC, and bootstrap intervals.
  The final gap estimates are not present in the supplied information.
  \item \textbf{Audit age removal and mitigation.} M1, M0-RW, and M1-RW are
  specified, including proxy ablation and performance-cost measurements.
  Their numerical outcomes require the final experiment output.
\end{enumerate}

\section{Future Work}

\begin{enumerate}
  \item \textbf{Complete repeated evaluation:} run the four model conditions
  on the primary and ten additional splits and populate the result tables with
  paired bootstrap intervals.
  \item \textbf{Use broader fairness criteria:} add equalized odds,
  calibration, and subgroup analyses for narrower age bands, while reporting
  the trade-offs caused by incompatible criteria.
  \item \textbf{Validate on real lending data:} evaluate on a dataset with
  documented collection procedures and reject-inference treatment, subject to
  privacy and governance requirements.
  \item \textbf{Study causal proxy mechanisms:} combine observational SHAP
  evidence with domain review and causal analysis before removing a feature
  that may also measure legitimate repayment capacity.
\end{enumerate}

The central lesson is that explainability and fairness should be evaluated
together with data quality. Removing a sensitive column is not sufficient when
other variables encode related information, and a high predictive score does
not by itself establish equitable treatment.
